\section{Represented real bodies and the strict volume premise}\label{sec:real}

\subsection{Positive lower volume as evidence}

A useful constructive convention defines the strict lower
volume assertion by its finite witnesses:
\begin{equation}\label{eq:lower-content}
 \Vol_-(D)>q
 \quad:\Longleftrightarrow\quad
 \begin{array}{l}
 \text{there is a finite family of rational boxes in \(D\),}\\[-.1em]
 \text{with disjoint interiors and total content greater than \(q\).}
 \end{array}
\end{equation}
For the arithmetic proof, containment refers to the
represented membership predicate on rational points.
For a real interpretation one uses actual certified
inclusion of the corresponding boxes.
This is a positive assertion: its witness contains
endpoints and inclusion evidence.

With this representation the strict theorem says exactly
\[
 \Sym(P)\wedge\Mid(P)\wedge\Vol_-(P/2)>1
 \ \Longrightarrow\
 \exists z\in\Z^n\setminus\{0\}\ (z\in P).
\]
It is the formula scheme of \cref{thm:main}, obtained by
existentially eliminating the supplied box certificate.
No supremum object or real-number completeness axiom
is needed.

For ordinary bounded convex bodies with nonempty interior,
lower box content agrees with the usual volume. The
relevant direction is visible in the shrink estimate:
an inradius absorbs the mesh error, and the content of
the shrink approaches that of the body. For rational
polytopes the preceding section has already replaced
this limiting sentence by one explicit mesh and finite
rational volume arithmetic.

For a different effective presentation of a real body,
the representation theorem connecting its volume to
\cref{eq:lower-content} must be proved. It is not
automatic for an opaque description of a measurable set.
For example, suppose a fixed approximation procedure
supplies certified inner unions \(U_t\) and a known
primitive recursive error bound. Given a rational
\(V_0>1\) below \(\Vol(D)\), a supplied or explicitly
computed \(t\) making the error smaller than
\((V_0-1)/2\) yields a union of content greater than
\(1\), hence a certificate. Merely knowing that some
unknown mesh eventually works does not give a
primitive recursive search bound.

An unbounded body is handled by the same certificate:
its positive lower volume witness is already a finite
bounded family of boxes. If one instead starts from
bounded intersections, the radius where the strict
threshold is exceeded must be part of the positive
evidence or supplied by a proved primitive recursive
bound.

\subsection{Real lattice bases}

For a real lattice \(L=B\Z^n\), the numerical output
remains the integer coefficient vector \(z\).
To interpret it as a nonzero real vector \(Bz\), use
constructive nonsingularity data: an explicitly
represented inverse, or a rational lower bound \(d>0\)
with \(|\det B|\geq d\), together with the arithmetic
of the fixed regular names. The bare negative
assertion \(\det B\neq0\) must not be silently treated
as such positive apartness data.

The predicate \(\Phi(q)\equiv Bq\in K\) can then be
expressed arithmetically for the fixed representation.
The box-collision proof is unchanged. It reads only the
finite coordinate certificate and uses \(\Phi\) in
its last three logical steps.

One may take \cref{eq:lower-content} in normalized
coordinates as the positive hypothesis corresponding
to \(\Vol(K)>2^n|\det B|\).
Alternatively, a separately proved volume-change rule
and positive inner approximation theorem connect it
with an existing volume representation. This makes
explicit where geometry is encoded, instead of
adding unrepresented real sets to a first-order
arithmetic theorem.

\section{Why unrestricted real equality is different}\label{sec:obstruction}

The rational equality argument depends on a discrete
integer-facet gap. It does not extend to arbitrary real
widths. The following complete reduction shows the
issue already for rectangles in dimension two.

\begin{proposition}[Equality gives weak real sign choice]\label{prop:llpo}
If the compact exact-threshold theorem produces a nonzero
integer point for every represented real symmetric
rectangle of area \(4\), then for every represented real
\(t\) with \(|t|\leq1/2\) one can produce a branch of
\[
                              t\geq0\ \vee\ t\leq0.
\]
\end{proposition}

\begin{proof}
Use
\[
 K_t=[-(1+t),1+t]\times
             [-(1+t)^{-1},(1+t)^{-1}].
\]
Since \(1+t\geq1/2\), the reciprocal is constructively
defined. The rectangle is bounded, closed, symmetric,
convex, and has area exactly \(4\).
Let \((u,v)\) be a nonzero integer point in it.
Integer equality is decidable.

If \(u\neq0\), then \(1\leq|u|\leq1+t\), so \(t\geq0\).
If \(u=0\), then \(v\neq0\) and
\(1\leq|v|\leq(1+t)^{-1}\). Multiplication by the
positive number \(1+t\) gives \(t\leq0\).
\end{proof}

Here is the direct LLPO reduction. Let \((\beta_k)\)
be a binary sequence with at most one term equal to
\(1\), and form the regular rational name
\[
                  t_s=\sum_{k<s}(-1)^k\beta_k2^{-(k+3)}.
\]
For \(v\geq s\),
\[
 |t_v-t_s|\leq\sum_{k\geq s}2^{-(k+3)}=2^{-(s+2)}.
\]
The modulus is explicit and \(|t|\leq1/4\).
If \(t\geq0\), no odd-indexed bit can equal \(1\):
such a bit would force all the other bits to be zero,
making \(t\) strictly negative. Hence all odd bits
are zero. Similarly \(t\leq0\) implies all even bits
are zero. The resulting disjunction is the lesser
limited principle of omniscience, LLPO.
The equivalence of weak real sign choice with LLPO
is discussed by Bridges and Richman
\cite{bridges-richman}.

This obstructs an unrestricted constructive
real-body extension. It does not contradict the finite
strict theorem or the rational equality theorem.
We do not claim to supply a Lean formalization of
LLPO independence from a particular coding of \(\ipra\).
The reduction is explicit, and the computational
obstruction can already be made first-order in one
fixed primitive recursive family.

\subsection{A uniform recursive selector is impossible}

Fix an effective enumeration \(\varphi_e\) of partial
recursive functions. The disjoint computably enumerable
sets
\[
 A=\{e:\varphi_e(e)\downarrow=0\},\qquad
 C=\{e:\varphi_e(e)\downarrow=1\}
\]
have no total recursive binary separator that is \(0\)
on \(A\) and \(1\) on \(C\). If \(f\) were such a
separator, let \(e\) be an index of the total function
\(x\mapsto1-f(x)\). If \(f(e)=0\), then \(e\in C\);
if \(f(e)=1\), then \(e\in A\). Each case contradicts
the required value of \(f(e)\).

For the parameter \(e\), define a rational approximation
\(t_e(s)\) by simulation for \(s\) steps. If
\(\varphi_e(e)\) first halts at step \(j\leq s\) with
output \(0\), use \(2^{-(j+3)}\); if it first halts
with output \(1\), use \(-2^{-(j+3)}\); otherwise use
\(0\). This is one fixed primitive recursive function
of \(e,s\), with an immediate geometric Cauchy modulus.
Its real value is positive on \(A\), negative on \(C\),
and zero otherwise.

A recursive selector for a nonzero integer point of
\(K_{t_e}\) would give the forbidden separator:
output \(0\) when the selected point has nonzero first
coordinate and \(1\) otherwise. For positive \(t_e\),
the second alternative is impossible; for negative
\(t_e\), the first is impossible.
This argument uses bounded universal Turing simulation,
not a total universal primitive recursive evaluator.

Under a realizability-sound first-order encoding, a
uniform constructive derivation producing these integer
witnesses would yield such a selector.
The package establishes the explicit reduction and
selector obstruction; it does not promote them to
an unimplemented syntactic independence theorem.
